\chapter{Pendahuluan}

\section{Apa itu PakAI?}

\textbf{PakAI} (dijalankan dengan perintah \cmd{pakai}) adalah alat bantu berbasis terminal
untuk memantau penggunaan langganan layanan AI secara terpadu. Alat ini menampilkan
sisa kuota dan biaya untuk beberapa penyedia AI sekaligus dalam satu tampilan yang
ringkas.

PakAI mendukung tiga penyedia layanan AI:

\begin{itemize}
  \item \textbf{Claude} (Anthropic): memantau persentase penggunaan pesan per 5 jam,
        mingguan, dan bulanan melalui OAuth API.
  \item \textbf{OpenAI Codex}: memantau kuota penggunaan per 5 jam dan mingguan
        melalui OAuth API.
  \item \textbf{OpenCode Go}: memantau biaya penggunaan per penyedia dari
        database lokal SQLite.
\end{itemize}

\section{Fitur Utama}

\begin{itemize}
  \item Pemantauan multi-penyedia dalam satu perintah.
  \item Tiga jendela waktu: 5 jam, mingguan, dan bulanan.
  \item Tampilan berwarna: hijau (aman), kuning (peringatan), merah (kritis).
  \item Empat mode tampilan: status CLI, tmux, Waybar, dan dashboard TUI.
  \item Daemon latar belakang dengan pembaruan otomatis dan API HTTP.
  \item Pengaturan melalui berkas konfigurasi TOML.
  \item Deteksi otomatis penyedia yang terinstal.
\end{itemize}

\section{Tampilan Aplikasi}

Gambar di bawah menampilkan output perintah \cmd{pakai status} yang menunjukkan
ringkasan penggunaan semua penyedia yang terdeteksi.

\begin{figure}[h]
  \centering
  \includegraphics[width=0.85\linewidth]{pakai-status}
  \caption{Output \texttt{pakai status}: ringkasan penggunaan semua penyedia}
\end{figure}

\section{Cara Kerja}

PakAI berjalan sebagai \textit{background daemon} yang secara berkala mengambil data
penggunaan dari setiap penyedia. Data ini kemudian ditampilkan melalui berbagai
antarmuka:

\begin{enumerate}
  \item \textbf{Daemon} berjalan di port 7731 dan menyimpan data terbaru.
  \item Perintah \cmd{pakai status}, \cmd{pakai tmux}, dan \cmd{pakai waybar}
        mengambil data dari daemon melalui HTTP.
  \item \cmd{pakai dashboard} menampilkan antarmuka TUI interaktif dengan
        pembaruan langsung (\textit{live update}).
\end{enumerate}
